\section{Lean Formalization Status}\label{sec:6}

The Lean formalization records the machine-checked definitions,
theorems, witnesses, and proof chain of the SAGA mathematics. This
section reports the formalization status at release time at declaration
granularity. The primary evidence for the status consists of the Lean
sources fixed by the release tag; the kernel axiom audit, which checks
every registered declaration against the allowlist of standard mathlib
axioms --- \code{propext}, \code{Classical.choice}, \code{Quot.sound};
and the focused check of the release CI (the organization of the
sources and the audit mechanism are in Appendix~\ref{app:B}). Every
declaration listed in the table of this section is registered in this
audit.

\subsection{Scope of the formalization}\label{sec:61}

The SAGA theorem chain is formalized by fixing the input structure of
the theorem of Section~\ref{sec:5} (inputs 1--8 of
Theorem~\ref{thm:central}) directly as Lean structures. Accordingly,
what each declaration proves is not the theorem of this paper as such,
but the conclusion over the selected inputs fixed as Lean structures
(the correspondence with the paper claims and the assumptions of each
row are fixed by the table of \S\ref{sec:62}). The building blocks are
as follows (the correspondence with source files is in
Appendix~\ref{app:B}):
\begin{itemize}
\item the monomorphic ordered cover, the intersection diagram, the
three-term \v{C}ech complex, and the cover-relative $H^1$;
\item the semantic repair presentation and the presheaf construction
of $\Msem$;
\item the affine semantic repair system, the selected atlases on both
sides, and the residuals;
\item the derivation of soundness, presentation exactness, the
coefficient isomorphism $\Phi$, and the finite counterexamples showing
the independence of the three conditions;
\item the equation-side realization / production (the connection to
the $\QE$ generated by the equation system);
\item the cochain comparison $\kappa$, the $H^1$ isomorphism, and the
residual correspondence;
\item true sheaf descent, grounded global gluing, and the final
bundling of the central theorem;
\item the finite witnesses (the nonzero 4-cycle circle witness and the
zero-class descent witness);
\item the site realization over the finite-meet poset model of
\S\ref{sec:33} and the comparison bridge to the ordered \v{C}ech
complex of \S\ref{sec:41} (including the formalized counterpart of
Lemma~5.2A).
\end{itemize}

The docstrings of the Lean sources use labels from a numbering series
distinct from the theorem numbers of this paper (\code{X.Theorem 1.1}
and so on). The correspondence between declarations and labels is also
fixed in Appendix~\ref{app:B}.

\subsection{Status table}\label{sec:62}

Every status in Table~\ref{tab:status} is \code{proved} (proved in Lean).

\begin{table}[htbp]
\centering
\scriptsize
\begin{tabular}{@{}p{0.27\textwidth}p{0.30\textwidth}p{0.06\textwidth}p{0.28\textwidth}@{}}
\toprule
Paper claim & Lean declaration & Status & Assumptions \\
\midrule
Conclusion bundle of Theorem~\ref{thm:central} (residual correspondence, zero/nonzero
equivalence, grounded gluing) &
\code{SagaEquationPacket.}\allowbreak\code{sagaCentralTheorem} & proved &
selected packet (\S\ref{sec:61}), the two completeness conditions
(input~4), \code{Fintype} of the cover index set. The gluing clause is
additionally conditional on the cover belonging to the topology and the
true sheaf condition \\
\addlinespace
Theorem~\ref{thm:central}(i): derivation of repair-relation soundness &
\code{PrimaryStateCorrespondence.}\allowbreak\code{relationSound\_}\allowbreak\code{of\_stateCorrespondence} &
proved & local-state correspondence (inputs 5--7) \\
\addlinespace
Theorem~\ref{thm:central}(i): coefficient isomorphism $\Phi$ &
\code{SagaEquationPacket.phiEquiv} (generic version:
\code{PrimaryCoefficientCorrespondence.}\allowbreak\code{phiEquiv}) & proved &
soundness and the two completeness conditions \\
\addlinespace
Theorem~\ref{thm:central}(i): counterexamples removing each of the three conditions &
\code{ExactnessFixtures.}\allowbreak\code{soundness\_failure} /
\code{completeness\_failure} / \code{generation\_failure} & proved &
none (each fixture is a finite counterexample stating the failure of
the targeted condition; simultaneous validity of the remaining two
conditions is not part of the statement) \\
\addlinespace
Theorem~\ref{thm:central}(ii): cochain commutation $\kappa\delta=\delta\kappa$ &
\code{kappa1\_delta0}, \code{kappa2\_delta1} & proved &
restriction-natural isomorphism of the coefficient families \\
\addlinespace
Theorem~\ref{thm:central}(ii): $H^1$ isomorphism $\kappa_*$ &
\code{SagaEquationPacket.}\allowbreak\code{kappaStarAddEquiv} ($\simeq_+$,
additive equivalence, on the packet; generic version:
\code{kappaH1AddEquiv}) & proved & the two
completeness conditions \\
\addlinespace
Theorem~\ref{thm:central}(ii): residual correspondence $\kappa_*([\rsem])=[\rE]$ &
\code{SagaEquationPacket.}\allowbreak\code{residual\_correspondence\_class};
aligned-atlas version \code{betaAligned\_residual} & proved & the two
completeness conditions \\
\addlinespace
Theorem~\ref{thm:gluing} = Theorem~\ref{thm:central}(iii): grounded global gluing &
\code{SagaEquationPacket.}\allowbreak\code{globalRepair\_nonempty\_iff}
($\Psem$-side equivalence); \code{sagaGroundedGluing} (integrated up to
the equation-side equivalence) & proved & true sheaf condition, cover
belonging to the topology; the equation-side equivalence additionally
requires the two completeness conditions \\
\addlinespace
Lemma~5.2A (ordered matching completion) &
\code{SiteStateData.matchingFamily\_iff} & proved & thin context
category (\S\ref{sec:57}, Appendix~\ref{app:A4}); subsingleton
assumption for states on the omitted pairs (the formalized counterpart
of empty-overlap normalization) \\
\addlinespace
Example~\ref{ex:circle} (4-cycle circle witness, transfer of the nonzero class) &
\code{CircleWitness.}\allowbreak\code{semanticResidualClass\_ne\_zero},
\code{circle\_nonzero\_class\_transfer} & proved & none (closed
verification on the concrete 4-cycle model) \\
\addlinespace
Zero-class witness (nonempty firing of Theorem~\ref{thm:gluing}) &
\code{DescentWitness.descentTrueSheaf},
\code{descent\_sagaGroundedGluing} & proved & none (closed verification
on a concrete model) \\
\bottomrule
\end{tabular}
\caption{Formalization status at declaration granularity.}
\label{tab:status}
\end{table}

All rows of the table are registered in the kernel axiom audit, and the
kernel axioms are limited to the standard mathlib axioms. No row
contains \code{sorry} or additional axioms.

What remains with respect to the mathematics of this paper
(Sections~\ref{sec:3}--\ref{sec:5}) --- to be distinguished from the
research outlook of Section~\ref{sec:9} --- is the following.
\begin{itemize}
\item The gluing argument of \S\ref{sec:56} over a general (non-thin)
context category, in the form that explicitly consumes monomorphisms
and the universal property of pullbacks: \code{unported} (the paper
proof exists but has not been ported to Lean; \S\ref{sec:57}).
\item Identification with sheaf cohomology via refinement invariance /
Leray-type acyclicity, nonabelian $H^1$, gerbes, stack descent: not
claimed by this paper, hence not among the proof obligations
(\S\ref{sec:57}).
\item Transporting the measurement runs of Section~\ref{sec:7} into
Lean (generating instantiations of Theorem~\ref{thm:central} from
measured packets): not done. The finite instantiation of
Theorem~\ref{thm:central} is carried by the witness rows of this table
(Example~\ref{ex:circle}), and the packets of Section~\ref{sec:7} carry
the finite checks shown by the condition matrix of Appendix~\ref{app:C3}. The
division of labor between the two is stated in \S\ref{sec:75}.
\end{itemize}

The release snapshot of this section is the repository state at the
release tag \code{saga-paper-v1.0.0}. The full Lean build, including
the axiom audit, runs in continuous integration on the tagged commit;
the commit hash and the reference to that CI run are recorded in
\code{MANIFEST.json} of the deposit bundle, which shares the release
identity of this paper (Appendix~\ref{app:C4}).
